\chapter{Percutaneous coronary intervention}
\index{Percutaneous coronary intervention}

\medskip
\noindent\textit{Educational information; seek professional advice for individual care.}

\section*{What it is}
Percutaneous coronary intervention (PCI), also called coronary angioplasty, uses a catheter to widen a narrowed or blocked coronary artery. A balloon may open the narrowing and a stent may hold the artery open. PCI can be planned for symptoms or performed urgently during a heart attack.

\section*{Before and during the procedure}
The cardiology team reviews symptoms, heart tests, kidney function, allergies, bleeding risk, pregnancy where relevant, and medicines. The catheter usually enters through an artery in the wrist or groin. Contrast imaging shows the blockage and guides treatment. Alternatives and the possibility of bypass surgery should be discussed when relevant.

\section*{After PCI}
Follow instructions about the access site, activity, hydration, and follow-up. Antiplatelet medicines are usually prescribed to prevent a clot forming in the stent. Do not stop or miss these medicines without urgent advice from the cardiology team, even if you feel well. Manage smoking, blood pressure, cholesterol, diabetes, activity, and cardiac rehabilitation as advised.

\section*{Risks}
Possible risks include bleeding, bruising, allergic or kidney effects from contrast, artery damage, abnormal rhythm, heart attack, stroke, stent clot, or restenosis. Report access-site swelling or bleeding and any new symptoms promptly.

\section*{When to seek emergency care}
Seek emergency help for chest pressure or pain, severe breathlessness, collapse, new neurological symptoms, or heavy bleeding. After discharge, urgent assessment is also needed for rapidly enlarging swelling at the access site, a cold or numb limb, or recurrent symptoms like the original angina.

\section*{Sources}
\begin{itemize}
\item American Heart Association: \url{https://www.heart.org/en/health-topics/heart-attack/treatment-of-a-heart-attack/stent}
\end{itemize}
